\begin{rubric}{Employment History}
\entry*[Jan 2026 - present]%
    \textbf{\href{https://eig.org/about-us/executive-team-staff/ben-glasner/}{Senior Economist}}. \href{https://eig.org/}{Economic Innovation Group (EIG)}, Washington, D.C.
    \par Set research direction: generate and prioritize projects against EIG's strategic priorities, and track economic, academic, and policy developments across labor markets, place-based policy, and retirement security.
    \par Lead research end to end: frame policy questions, select data, design and defend identification strategies, build microsimulation and fiscal-cost models, and write publication-ready analysis. Methods: applied microeconomics, modern difference-in-differences, geospatial analysis, and large administrative and survey datasets.
    \par Build research infrastructure: set EIG's standards for reproducible research (GitHub collaboration, code review, and replication packages) and for AI-assisted empirical work (AI coding agents used under verification and disclosure standards).
    \par Represent EIG publicly, translating technical work for policymakers, journalists, and broadcast audiences.
    \par Mentor and develop junior staff on coding, data practices, and quality norms; support fundraising and partnerships with funders, contractors, and vendors.
\entry*[Sep 2024 - Jan 2026]%
    \textbf{Economist}. \href{https://eig.org/}{Economic Innovation Group (EIG)}, Washington, D.C.
    \par Led quasi-experimental policy evaluations with large administrative and survey datasets, including advanced statistical and spatial analysis at the county and census-tract levels.
    \par Developed evidence-based policy proposals, from literature review and data analysis through program design.
    \par Translated empirical results into policy writing, data visualization, and stakeholder outreach.
\entry*[Jan 2023 - Sep 2024]%
    \textbf{Associate Economist}. \href{https://eig.org/}{Economic Innovation Group (EIG)}, Washington, D.C.
    \par Applied empirical methods to research on labor markets, business formation, remote work, and place-based policy.
    \par Built data pipelines and supported report production, from data engineering through publication and policy communication; synthesized research literatures for policy audiences.
\entry*[2021 - 2022]%
    \textbf{\href{https://povertycenter.columbia.edu/people}{Postdoctoral Research Scientist}}. \href{https://povertycenter.columbia.edu/people}{Center on Poverty and Social Policy, Columbia University}, New York, NY
    \par Empirical analyses of major U.S. social policies and reforms. Work on intergenerational poverty, material hardship, and employment impacts of the Child Tax Credit.
\entry*[2020 - 2021]%
    \textbf{Dean's Research Assistant}. University of Washington, Seattle, WA
    \par Supported faculty transition to technology-enhanced learning. Conducted DEI and course design assessments; developed action plans informed by interviews and surveys.
\entry*[2020]%
    \textbf{Predoctoral Instructor}. University of Washington, Seattle, WA
    \par Primary instructor for a section of the core MPA microeconomics and public policy course.
\entry*[2018]%
    \textbf{Short-Term Consultant}. World Bank Group, Seattle, WA
    \par Expenditure analysis across household surveys in Malawi, Tanzania, Ghana, Peru, and Cote d'Ivoire. Developed seasonality estimates using parametric, semi-parametric, and machine learning approaches.
\end{rubric}
